\section{Methodology}
\label{sec:method}

\subsection{Problem statement and nondimensionalization}

We seek a map from a surface descriptor $\mathbf{s}$ and a flow condition
$\mathbf{f}$ to the ONB wall superheat and heat flux,
\begin{equation}
(\mathbf{s},\mathbf{f}) \;\longmapsto\; (\dTonb,\ \qonb),
\end{equation}
where $\dTonb = T_w - T_{\mathrm{sat}}(P)$ and $\qonb$ is the heat flux at
incipience. Targets are nondimensionalised per sample using fluid properties
at the local saturation state: $\dTonb^{\*} = \dTonb/\dTonb_{\mathrm{ref}}$ with
a fixed reference $\dTonb_{\mathrm{ref}}=\SI{20}{K}$ (mid-range of the
physical $2$--\SI{30}{K} window for water), and $q^{\*}=\qonb/q_{\mathrm{ref}}$
with a capillary-scaled $q_{\mathrm{ref}}$. Fluid properties are evaluated with
CoolProp \cite{COOLPROP}. The axial coordinate $x^{\*}=x/L\in[0,1]$ locates the
ONB site; data reported without an axial location use $x^{\*}=0.5$.

\subsection{Architecture}
\label{sec:arch}

The network (Fig.~\ref{fig:arch}) has three parts.

\begin{figure}[t]
\centering
\includegraphics[width=\linewidth]{architecture.png}
\caption{Surface- and flow-conditioned PINN. A frozen surface encoder
(transfer-learned from the Phase~1 pool-boiling model) and a new flow encoder
produce latent vectors $\mathbf{z}_s,\mathbf{z}_f$ that condition a six-layer
backbone in the axial coordinate $x^{*}$ through feature-wise linear modulation
(FiLM). Three heads output the ONB wall superheat $\dTonb^{*}$, heat flux
$q^{*}$ and an auxiliary temperature field $T^{*}(x^{*})$. Training combines a
masked data loss with a flow-aware Hsu collocation constraint, a one-sided Hsu
self-consistency coupling and monotonicity penalties.}
\label{fig:arch}
\end{figure}

\paragraph{Surface encoder (transfer-learned, frozen).}
Surface morphology---arithmetic roughness $R_a$, static contact angle $\theta$
and a categorical surface/treatment label---is embedded by the encoder trained
in Phase~1 on pool-boiling data. The upstream feature-extraction layers are
loaded from the Phase~1 checkpoint and held frozen; only a thin projection to
the Phase~2 latent dimension is re-initialised (the latent widths differ,
$32\!\to\!16$). Freezing both preserves the pool-boiling surface knowledge and
prevents the small flow-boiling dataset from corrupting it.

\paragraph{Flow encoder (new).}
Flow conditions are encoded by an eleven-channel descriptor: $\log_{10}\!Re/5$,
$\log_{10}(G/G_{\mathrm{ref}})$, $\log_{10}(Bo\!\cdot\!10^4{+}1)$,
$\log_{10}(We{+}1)$, the subcooling ratio $\dTsub/T_{\mathrm{sat}}$ with a
provided/missing mask, $\log_{10}(\Dh{+}0.01)$, a sine--cosine encoding of
channel geometry (tube / rectangular / microchannel), and the
\emph{reduced pressure} $\log_{10}(P/P_{\mathrm{crit}})$ with a
provided/missing mask. The reduced-pressure channel is essential: without it
(Section~\ref{sec:patternC}) the network cannot represent the strong
$P\!\uparrow\!\Rightarrow\!\dTonb\!\downarrow$ dependence, because that
dependence enters property-based groups only weakly. Missing pressures
(literature reporting gaps) are imputed at atmospheric with the mask set to
zero, letting the network distinguish measured from imputed values.

\paragraph{Backbone and heads.}
The two latent vectors condition a six-layer multilayer-perceptron backbone in
$x^{\*}$ through feature-wise linear modulation (FiLM). Three positive
soft-plus heads output $\dTonb^{\*}$, $q^{\*}$ and an auxiliary temperature
field $T^{\*}(x^{\*})$. The model has $\sim$\,$9.9\times10^{4}$ parameters.

\subsection{Loss function}
\label{sec:loss}

The model is a \emph{physics-regularized regressor}: it fits the scalar ONB
data while soft physical constraints shape the solution. The total loss is
\begin{equation}
\mathcal{L} = w_q\mathcal{L}_q + w_{\Delta T}\mathcal{L}_{\Delta T}
+ w_{\mathrm{Hsu}}\mathcal{L}_{\mathrm{Hsu}}
+ w_{\mathrm{cp}}\mathcal{L}_{\mathrm{cp}}
+ \!\!\sum_{\xi\in\{Re,G,\dTsub\}}\!\! w_\xi\,\mathcal{L}^{\mathrm{mono}}_{\xi}.
\label{eq:loss}
\end{equation}
$\mathcal{L}_q,\mathcal{L}_{\Delta T}$ are masked errors on the available
targets (log-scaled for $q$). $\mathcal{L}_{\mathrm{Hsu}}$ is a flow-aware Hsu
constraint on collocation points: the Hsu discriminant uses a thermal-layer
thickness $\delta_t = C\,\Dh/\mathrm{Nu}$ with $\mathrm{Nu}$ from
Dittus--Boelter, so that nucleation feasibility responds to the flow.

\paragraph{One-sided Hsu coupling.}
$\mathcal{L}_{\mathrm{cp}}$ enforces self-consistency between the two heads
through the Hsu relation $\dTonb^{\*}=C_{\mathrm{Hsu}}\sqrt{q^{\*}}$. Crucially
this is applied as a \emph{one-sided lower bound},
\begin{equation}
\mathcal{L}_{\mathrm{cp}} = \big\langle\,
\mathrm{ReLU}\!\big(C_{\mathrm{Hsu}}\sqrt{q^{\*}}-\dTonb^{\*}\big)^2\big\rangle,
\label{eq:onesided}
\end{equation}
penalising predictions \emph{below} the saturated-pool Hsu superheat but not
above it. The rationale is physical: bulk subcooling and forced convection
\emph{raise} the wall superheat required for incipience, so the
saturated-pool value is a lower bound, $\dTonb \ge \dTonb^{\mathrm{Hsu,sat}}$.
An equality coupling (the naive choice) instead drags the high-subcooling,
low-heat-flux predictions toward the near-zero saturated value and biases them
low; the one-sided form removes that spurious pull
(Section~\ref{sec:patternB}). Figure~\ref{fig:concept} illustrates both the
identification of ONB from a boiling curve and the lower-bound interpretation of
the Hsu criterion.

\begin{figure}[t]
\centering
\includegraphics[width=\linewidth]{fig_onb_concept.png}
\caption{(a)~ONB is identified as the departure of the wall temperature from its
single-phase line; the excess over saturation is $\dTonb$. (b)~The
saturated-pool Hsu superheat $C\sqrt{q^{*}}$ is a \emph{lower bound}: subcooling
and convection place the measured subcooled-flow ONB above it, motivating the
one-sided coupling of Eq.~\eqref{eq:onesided}.}
\label{fig:concept}
\end{figure}

\paragraph{Monotonicity.}
$\mathcal{L}^{\mathrm{mono}}_{\xi}$ penalise violations of the expected sign of
$\partial\dTonb/\partial\xi$ along one-dimensional sweeps of Reynolds number,
mass flux and subcooling, evaluated on synthetic collocation points.

\paragraph{On full PDE residuals.}
We deliberately do \emph{not} impose Navier--Stokes (continuity, momentum)
residuals. The model outputs scalar incipience quantities, not resolved
velocity and pressure fields, so those residuals are undefined here. We
explored adding a one-dimensional advection-dominated energy residual on the
auxiliary $T^{\*}(x^{\*})$ head; it destabilised training (its second-order
autograd gradients dominate the gradient-clip budget and starve the data loss)
while regularising only an unsupervised field, and was not adopted. Full
field-resolving treatment is left to future work; the present
physics-regularized formulation already enforces the relevant scalar trends
(Section~\ref{sec:physcon}).

\subsection{Training}

Parameters are optimised with a two-stage schedule (Adam with cosine decay,
then L-BFGS polishing), the surface encoder frozen throughout. Hsu and
monotonicity collocation batches are resampled each epoch. Unless noted, the
data are split 70/15/15 stratified by source study; for the deep ensemble
(Section~\ref{sec:uq}) the split is held fixed while the initialisation seed is
varied.
